\documentclass[12pt,a4paper]{article}
\usepackage[margin=25mm, headheight=15pt, headsep=10pt]{geometry}
\usepackage{fontspec}
\usepackage[table]{xcolor} % Note: table option is required for rowcolors
\usepackage{titlesec}
\usepackage{tocloft}
\usepackage{fancyhdr}
\usepackage{booktabs}
\usepackage{tcolorbox}
\tcbuselibrary{skins, breakable}
\usepackage[colorlinks=true, linkcolor=emerald, urlcolor=emerald, bookmarks=true]{hyperref}

\definecolor{emerald}{RGB}{0,102,51}
\definecolor{darkred}{RGB}{139,0,0}
\definecolor{lightgray}{RGB}{245,245,245}

\usepackage[nil]{babel}
\babelprovide[import, main]{bengali}
\babelprovide[import]{arabic}
\babelfont{rm}[Renderer=HarfBuzz,Scale=1.0]{Noto Serif Bengali}
\babelfont{sf}[Renderer=HarfBuzz]{Noto Sans Bengali}
\babelfont{tt}[Renderer=HarfBuzz]{Noto Sans Bengali}
\babelfont[arabic]{rm}[Renderer=HarfBuzz,Scale=1.1]{Noto Naskh Arabic}

% Section styling
\titleformat{\section}{\Large\bfseries\color{emerald}}{\thesection.}{0.5em}{}
\titleformat{\subsection}{\large\bfseries\color{emerald!85!black}}{\thesubsection.}{0.5em}{}
\titleformat{\subsubsection}{\normalsize\bfseries\color{emerald!70!black}}{\thesubsubsection.}{0.5em}{}
\titlespacing*{\section}{0pt}{18pt}{8pt}

% Fancy Header/Footer setup
\pagestyle{fancy}
\fancyhf{} % clear all header and footer fields
\fancyhead[L]{\color{emerald}\small\textbf{\nouppercase{\leftmark}}}
\fancyfoot[L]{\scriptsize\color{black!50}{\lat{AI}}-সংকলিত, আলেম-যাচাইকৃত নয়}
\fancyfoot[C]{\color{emerald!70!black}\thepage}
\renewcommand{\headrulewidth}{0.4pt}
\renewcommand{\headrule}{\hbox to\headwidth{\color{emerald}\leaders\hrule height \headrulewidth\hfill}}
\renewcommand{\footrulewidth}{0pt}

% Custom boxes
% 1. Ayat/Hadith Box
\newtcolorbox{ayatbox}[1][]{
    enhanced, breakable, 
    colback=emerald!5, 
    colframe=emerald, 
    boxrule=0.5pt, 
    arc=3pt, 
    left=10pt, right=10pt, top=10pt, bottom=10pt,
    #1
}

% 2. Quote/Translation Box
\newtcolorbox{quotebox}[1][]{
    enhanced, breakable,
    frame hidden,
    colback=lightgray,
    borderline west={3pt}{0pt}{emerald},
    left=15pt, right=10pt, top=10pt, bottom=10pt,
    sharp corners,
    #1
}

% 3. AI-generated notice box
\newtcolorbox{warnbox}[1][]{
    enhanced, breakable,
    colback=red!4,
    colframe=darkred,
    boxrule=0.8pt,
    arc=3pt,
    left=10pt, right=10pt, top=10pt, bottom=10pt,
    #1
}

% Commands
\newcommand{\arab}[1]{\foreignlanguage{arabic}{#1}}
\newcommand{\kib}[1]{\textcolor{emerald}{\textbf{#1}}}

% Latin (English) fallback: Noto Serif Bengali has NO Latin glyphs
\newfontfamily\latfont[Renderer=HarfBuzz]{Noto Serif}
\newcommand{\lat}[1]{{\latfont #1}}

\setlength{\parskip}{5pt}
\setlength{\parindent}{0pt}

% Fix for missing bullet in Noto Serif Bengali
\renewcommand{\labelitemi}{$\bullet$}



\begin{document}
\begin{center}{\Large\bfseries\color{emerald} হাদিস ও আসার — অহংকার (কিবর) ও আত্মদাম্ভিকতা (উজব)}\end{center}
\vspace{1em}
\section*{হাদিস ও আসার — অহংকার (কিবর) ও আত্মদাম্ভিকতা (উজব)}

\begin{warnbox}
 \textbf{গুরুত্বপূর্ণ নোটিশ (\lat{Important} \lat{Notice}):} এই কন্টেন্টটি \lat{AI} (কৃত্রিম বুদ্ধিমত্তা) দিয়ে প্রস্তুত ও সংকলিত। \lat{AI} কঠোর সূত্র-মান নীতি মেনে শুধু নির্ভরযোগ্য উৎস থেকে তথ্য সংগ্রহ করেছে — কেবল সহীহ/হাসান হাদিস ও সহীহ সনদের আসার রাখা হয়েছে, দুর্বল সনদ বাদ দেওয়া হয়েছে। তবে এটি কোনো আলেম বা মুহাদ্দিস কর্তৃক যাচাইকৃত নয়।
\end{warnbox}


\begin{quote}
\textbf{পড়ার নিয়ম:} এই ফাইলের প্রতিটি হাদিসের সাথে থাকবে — \textbf{সূত্র কার্ড} (সনদের মান) $\rightarrow$ \textbf{পটভূমি} $\rightarrow$ \textbf{পূর্ণ বর্ণনা} $\rightarrow$ \textbf{বিশ্লেষণ} $\rightarrow$ \textbf{শব্দের ব্যাখ্যা} $\rightarrow$ \textbf{গুরুত্বপূর্ণ শিক্ষা}। \\
\textbf{মূল নীতিমালা:} শুধুমাত্র সহীহ/হাসান হাদিস এবং সহীহ সনদের আসার (আথার) গ্রহণ করা হয়েছে। দুর্বল সনদ সম্পূর্ণ বাদ দেওয়া হয়েছে — লেবেল দিয়ে রাখাও হয়নি।
\end{quote}


\medskip\hrule\medskip

\subsection*{পার্ট ১ — কিবরের মৌলিক নিষেধ ও সংজ্ঞা}

\subsubsection*{হাদিস ১ — সরিষার দানাপরিমাণ কিবর + কিবরের সংজ্ঞা (সহীহ মুসলিম ৯১এ)}

\#\#\#\# ১. সূত্র কার্ড

\begin{table}[h]\centering\small
\begin{tabular}{ll}
বিষয় & বিবরণ \\
\hline
\textbf{বর্ণনাকারী} & আবদুল্লাহ ইবনে মাসউদ (রাঃ) \\
\textbf{গ্রন্থ} & সহীহ মুসলিম ৯১ (এ); সহীহ বুখারি ৩০/৬০৫০; সুনানে তিরমিজি ১৯৯৯; সুনানে আবু দাউদ ৪০৯১ \\
\textbf{অধ্যায়} & কিতাবুল ইমান — বাবু তাহরিমিল কিবর (কিবর হারাম হওয়ার অধ্যায়) \\
\textbf{সনদের মান} & \textbf{সহীহ} (বুখারি-মুসলিম) \\
\hline
\end{tabular}
\end{table}

\#\#\#\# ২. পটভূমি (\lat{Context})

এটি কিবর-উজব বিষয়ের সবচেয়ে গুরুত্বপূর্ণ হাদিস। ইবনে মাসউদ (রাঃ) — কুরআনের বিশিষ্ট কারী ও নবীজির ঘরের সদস্য — এই হাদিস বর্ণনা করেছেন। প্রশ্নটি ছিল অত্যন্ত \lat{practical}: মানুষ কি সুন্দর পোশাক-জুতা পছন্দ করলে সেটিও কি কিবর? নবীজি (সা.) এই প্রশ্নের মাধ্যমে কিবরের \textbf{\lat{clear} \lat{definition} (স্পষ্ট সংজ্ঞা)} দিয়েছেন — যা সমগ্র আলেমদের কাছে মানদণ্ড।

\#\#\#\# ৩. পূর্ণ বর্ণনা (বাংলা, \lat{pure})

\begin{quote}
আবদুল্লাহ বিন মাসউদ (রাঃ) বর্ণনা করেন, রাসূলুল্লাহ (সাল্লাল্লাহু আলাইহি ওয়া সাল্লাম) বলেছেন: \\
\textbf{\lat{“}যার হৃদয়ে সরিষার দানার ওজনের কিবর থাকবে, সে জান্নাতে প্রবেশ করবে না।\lat{”}} \\
একজন ব্যক্তি জিজ্ঞেস করলেন: \textbf{\lat{“}হে আল্লাহর রাসূল! মানুষ ভালো পোশাক ও জুতা পছন্দ করে। কি এটাও কিবর?\lat{”}} \\
নবীজি (সা.) উত্তর দিলেন: \textbf{\lat{“}আল্লাহ সুন্দর, তিনি সৌন্দর্য ভালোবাসেন। কিবর হলো সত্যকে তুচ্ছ করা ও মানুষকে হেয় করা।\lat{”}}
\end{quote}

\#\#\#\# ৪. ধাপে ধাপে বিশ্লেষণ

\textbf{ধাপ ১ — সতর্কবার্তা:} \lat{“}যার হৃদয়ে সরিষার দানার কিবর\lat{”} — অতি সামান্য পরিমাণ অহংকারও জান্নাতের পথে বাধা। এটি \textbf{\lat{magnitude} (পরিমাণ)} নয়, উপস্থিতির প্রশ্ন।

\textbf{ধাপ ২ — প্রশ্ন:} সাহাবি জিজ্ঞেস করলেন — ভালো পোশাক ভালোবাসা কি কিবর? অর্থাৎ কিবর কি বাহ্যিক সৌন্দর্যের সাথে জড়িত?

\textbf{ধাপ ৩ — উত্তর:} নবীজি দুটি জিনিস আলাদা করলেন — (ক) সৌন্দর্য পছন্দ করা বৈধ ও ভালো (আল্লাহ নিজেই সুন্দর), (খ) কিবর একটি ভেতরের অবস্থা — সত্যকে তুচ্ছ করা ও মানুষকে হেয় করা।

\textbf{ধাপ ৪ — সংজ্ঞা:} কিবর মানে দামি পোশাক নয়; কিবর মানে হককে মেনে না নেওয়া এবং মানুষকে ছোট করা।

\#\#\#\# ৫. শব্দের ব্যাখ্যা (\lat{Vocabulary})

\begin{table}[h]\centering\small
\begin{tabular}{ll}
শব্দ & অর্থ \\
\hline
\textbf{মিছকাল যাররাহ (\arab{مثقال} \arab{ذرة})} & সরিষার দানার ওজন — অতি সামান্য পরিমাণ \\
\textbf{কিবর (\arab{كبر})} & অহংকার — নিজেকে বড় ভাবা \\
\textbf{বাত্বারুল হক্ক (\arab{بطر} \arab{الحق})} & সত্যকে তুচ্ছ করা — হক জেনে মেনে না নেওয়া \\
\textbf{গামত্বুন নাস (\arab{غمط} \arab{الناس})} & মানুষকে হেয় করা — ছোট করে দেখা \\
\textbf{যামীল (\arab{جميل})} & সুন্দর — আল্লাহর গুণ \\
\hline
\end{tabular}
\end{table}

\#\#\#\# ৬. গুরুত্বপূর্ণ শিক্ষা (\lat{Key} \lat{Lessons})

\begin{enumerate}
\item \textbf{কিবরের সঠিক সংজ্ঞা:} দুটি জিনিস — (ক) সত্যকে তুচ্ছ করা, (খ) মানুষকে হেয় করা। যারা এই দুটি থেকে বাঁচে, তারা কিবর থেকে নিরাপদ।
\item \textbf{সৌন্দর্য পছন্দ করা জায়েজ:} ইসলাম দরিদ্রতা নয়; পরিচ্ছন্ন, সুন্দর পোশাক পছন্দ করা সুন্নাত — যতক্ষণ না তাতে অহংকার মেশে।
\item \textbf{সামান্য কিবরও বিপদ:} সরিষার দানার মতো ছোট কিবরও জান্নাতের পথে বাধা — এই সতর্কতা প্রতিটি মুমিনের \textbf{\lat{accountability} (দায়বদ্ধতা)}।
\item \textbf{অহংকার ভেতরের রোগ:} কিবর বাহ্যিক সম্পদে নয়, ভেতরের মনোভাবে।
\end{enumerate}
\#\#\#\# ৭. জীবনে প্রয়োগ (\lat{Practical} \lat{Application})

\begin{enumerate}
\item সত্য শুনে মেনে নেওয়ার \textbf{\lat{habit} (অভ্যাস)} গড়ুন — এমনকি সেটা নিজের বিরুদ্ধে হলেও।
\item মানুষকে ছোট করে কথা বলা, হাসি-ঠাট্টা — এড়িয়ে চলুন।
\item সুন্দর পোশাক পরুন, কিন্তু নিয়ত রাখুন — পরিচ্ছন্নতা ও আল্লাহর নেয়ামতের শুকরিয়া, অহংকার নয়।
\end{enumerate}
\medskip\hrule\medskip

\subsubsection*{হাদিস ২ — ঈমান ও কিবরের একই ওজন (সহীহ মুসলিম ৯১বি)}

\#\#\#\# ১. সূত্র কার্ড

\begin{table}[h]\centering\small
\begin{tabular}{ll}
বিষয় & বিবরণ \\
\hline
\textbf{বর্ণনাকারী} & আবদুল্লাহ ইবনে মাসউদ (রাঃ) \\
\textbf{গ্রন্থ} & সহীহ মুসলিম ৯১বি; সুনানে ইবনে মাজাহ ৫৯, ৪১৭৩; সুনানে তিরমিজি ১৯৯৮; মুসনাদ আহমদ ৩৯১৩ \\
\textbf{সনদের মান} & \textbf{সহীহ} \\
\hline
\end{tabular}
\end{table}

\#\#\#\# ২. পটভূমি (\lat{Context})

এটি একই মজলিসের ধারাবাহিক হাদিস। নবীজি (সা.) দুটি বিপরীত প্রতিশ্রুতি দিচ্ছেন — ঈমানের ওজন আছে তাকে জাহান্নাম থেকে বাঁচায়, আর কিবরের ওজন আছে তাকে জান্নাত থেকে বঞ্চিত করে। এই \textbf{\lat{balance} (তুলনা)} অত্যন্ত ভয়াবহ — কারণ দুটি জিনিসই মানুষের হৃদয়ে একসাথে থাকতে পারে।

\#\#\#\# ৩. পূর্ণ বর্ণনা (বাংলা, \lat{pure})

\begin{quote}
\textbf{\lat{“}যার হৃদয়ে সরিষার দানার ওজনে ঈমান থাকবে, সে জাহান্নামে প্রবেশ করবে না; আর যার হৃদয়ে সরিষার দানার ওজনে কিবর থাকবে, সে জান্নাতে প্রবেশ করবে না।\lat{”}}
\end{quote}

\#\#\#\# ৪. ধাপে ধাপে বিশ্লেষণ

\textbf{ধাপ ১ — ঈমানের প্রতিশ্রুতি:} সামান্য ঈমানও জাহান্নাম থেকে রক্ষা করে — আল্লাহর রহমতের প্রশস্ততা।

\textbf{ধাপ ২ — কিবরের প্রতিশ্রুতি:} সামান্য কিবরও জান্নাত থেকে বঞ্চিত করে — আল্লাহর সতর্কতার কঠোরতা।

\textbf{ধাপ ৩ — মাঝপথ:} এখানে 'খালিদ ফিল জাহান্নাম' বোঝানো হয়নি; ইমাম নববির ব্যাখ্যা — কিবরওয়ালা যতদিন কিবর থাকবে, ততদিন জান্নাতে প্রবেশ পাবে না; পরে শাস্তি ভোগ করে প্রবেশ করতে পারে। তবে এটি কোনো স্বস্তির কথা নয় — কারণ কিবর ছেড়ে যাওয়া কঠিন।

\#\#\#\# ৫. শব্দের ব্যাখ্যা (\lat{Vocabulary})

\begin{table}[h]\centering\small
\begin{tabular}{ll}
শব্দ & অর্থ \\
\hline
\textbf{যাররাহ (\arab{ذرة})} & ছোট পরিমাণ — সরিষার দানা \\
\textbf{কিবরিয়া (\arab{كبرياء})} & অহংকার — এখানে কিবরের চরম রূপ \\
\textbf{নূর / ঈমান} & ঈমান — প্রত্যয় \\
\hline
\end{tabular}
\end{table}

\#\#\#\# ৬. গুরুত্বপূর্ণ শিক্ষা (\lat{Key} \lat{Lessons})

\begin{enumerate}
\item \textbf{ঈমান ও কিবর এক হৃদয়ে প্রতিদ্বন্দ্বী:} মুমিনকে সতর্ক থাকতে হবে — যেন কিবর তার ঈমানের ওপর \textbf{\lat{overpower} (প্রভাব)} না ফেলে।
\item \textbf{সামান্য জিনিসের গুরুত্ব:} আল্লাহর কাছে পরিমাণ নয়, \textbf{\lat{quality} (গুণ)} গুরুত্বপূর্ণ — সামান্য ঈমানও মূল্যবান, সামান্য কিবরও মারাত্মক।
\item \textbf{আশা ও ভয়:} একই হাদিসে আশা (ঈমানের পুরস্কার) ও ভয় (কিবরের শাস্তি) — উভয়ই।
\end{enumerate}
\#\#\#\# ৭. জীবনে প্রয়োগ (\lat{Practical} \lat{Application})

\begin{enumerate}
\item নিজের হৃদয় পরীক্ষা করুন — কিবরের কোনো চিহ্ন আছে কি না।
\item ঈমান বাড়ানোর আমল (জিকির, নামাজ, কুরআন) নিয়মিত করুন — যেন ঈমানের ওজন বেশি থাকে।
\end{enumerate}
\medskip\hrule\medskip

\subsubsection*{হাদিস ৩ — হাদিস কুদসি: কিবর আল্লাহর চাদর (সহীহ মুসলিম ২৬২০)}

\#\#\#\# ১. সূত্র কার্ড

\begin{table}[h]\centering\small
\begin{tabular}{ll}
বিষয় & বিবরণ \\
\hline
\textbf{বর্ণনাকারী} & আবু সাঈদ ও আবু হুরাইরা (রাঃ) \\
\textbf{গ্রন্থ} & সহীহ মুসলিম ২৬২০; সুনানে আবু দাউদ ৪০৯০; সুনানে ইবনে মাজাহ ৪১৭৪ \\
\textbf{অধ্যায়} & কিতাবুল বিরর — বাবু তাহরিমিল কিবর \\
\textbf{সনদের মান} & \textbf{সহীহ} \\
\hline
\end{tabular}
\end{table}

\#\#\#\# ২. পটভূমি (\lat{Context})

এটি একটি \textbf{\lat{hadith} \lat{qudsi} (হাদিস কুদসি)} — যেখানে আল্লাহ নিজে কথা বলেছেন। অলংকারের ভাষায় আল্লাহ বলছেন — কিবর ও মহিমা আমার একান্ত \textbf{\lat{exclusive} (একান্ত)} গুণ; কেউ যদি সেটি ছিনিয়ে নিতে চায়, তবে সে আমার সাথে লড়াই করল — আর তার শাস্তি জাহান্নাম।

\#\#\#\# ৩. পূর্ণ বর্ণনা (বাংলা, \lat{pure})

\begin{quote}
আবু সাঈদ ও আবু হুরাইরা (রাঃ) থেকে: রাসূলুল্লাহ (সাল্লাল্লাহু আলাইহি ওয়া সাল্লাম) বললেন, আল্লাহু তাআলা বলেছেন: \\
\textbf{\lat{“}কিবর আমার চাদর, মহিমা আমার লুঙ্গি। যে আমার সাথে এর কোনো একটিতে প্রতিযোগিতা করবে, আমি তাকে জাহান্নামে নিক্ষেপ করব।\lat{”}}
\end{quote}

\#\#\#\# ৪. ধাপে ধাপে বিশ্লেষণ

\textbf{ধাপ ১ — রূপক:} চাদর ও লুঙ্গি — রাজার পোশাকের দুটি অংশ; এগুলো আল্লাহর \textbf{\lat{greatness} (মহত্ত্ব)}-এর প্রতীক।

\textbf{ধাপ ২ — প্রতিযোগিতা:} যে নিজেকে বড় ভাবে, মহিমার দাবি করে — সে আল্লাহর সাথে এই দুই গুণে \textbf{\lat{competition} (প্রতিযোগিতা)} করে।

\textbf{ধাপ ৩ — পরিণতি:} প্রতিযোগীর শাস্তি জাহান্নামে নিক্ষেপ — অহংকারের সবচেয়ে কঠিন শাস্তির ঘোষণা।

\#\#\#\# ৫. শব্দের ব্যাখ্যা (\lat{Vocabulary})

\begin{table}[h]\centering\small
\begin{tabular}{ll}
শব্দ & অর্থ \\
\hline
\textbf{কিবর (\arab{كبر})} & অহংকার — বড়ত্বের দাবি \\
\textbf{কিবরিয়া (\arab{كبرياء})} & মহিমা — আল্লাহর একান্ত গুণ \\
\textbf{আযামাহ (\arab{عظمة})} & মহত্ত্ব \\
\textbf{রিদা (\arab{رداء})} & চাদর \\
\textbf{ইযার (\arab{إزار})} & লুঙ্গি — কোমরের কাপড় \\
\textbf{নাযাআ} & প্রতিযোগিতা — ছিনিয়ে নেওয়ার চেষ্টা \\
\hline
\end{tabular}
\end{table}

\#\#\#\# ৬. গুরুত্বপূর্ণ শিক্ষা (\lat{Key} \lat{Lessons})

\begin{enumerate}
\item \textbf{অহংকার — আল্লাহর একান্ত ক্ষেত্রে আগ্রাসন:} নিজেকে বড় ভাবা মানে আল্লাহর মহত্ত্বের \textbf{\lat{territory} (অঞ্চল)}-এ ঢোকা — এটি চরম অন্যায়।
\item \textbf{বিনয় — আল্লাহর গুণের প্রতি সম্মান:} বিনয়ী ব্যক্তি আল্লাহর মহত্ত্ব স্বীকার করে — তাই সে আল্লাহর প্রিয়।
\item \textbf{জাহান্নামের সরাসরি শাস্তি:} এই হাদিসে অহংকারের পরিণতি সরাসরি বলা হয়েছে।
\end{enumerate}
\#\#\#\# ৭. জীবনে প্রয়োগ (\lat{Practical} \lat{Application})

\begin{enumerate}
\item নিজের বড়ত্বের দাবি ত্যাগ করুন — বড়ত্ব শুধু আল্লাহর।
\item অহংকার অনুভব করলে সাথে সাথে \lat{“}\arab{لا} \arab{حول} \arab{ولا} \arab{قوة} \arab{إلا} \arab{بالله}\lat{”} পড়ুন এবং বিনয়ের আমল করুন।
\end{enumerate}
\medskip\hrule\medskip

\subsection*{পার্ট ২ — পোশাক ও চলনে অহংকার}

\subsubsection*{হাদিস ৪ — পোশাক টেনে হাঁটা (সহীহ বুখারি ৫৭৮৪, মুসলিম ২০৮৫)}

\#\#\#\# ১. সূত্র কার্ড

\begin{table}[h]\centering\small
\begin{tabular}{ll}
বিষয় & বিবরণ \\
\hline
\textbf{বর্ণনাকারী} & আবু হুরাইরা (রাঃ) \\
\textbf{গ্রন্থ} & সহীহ বুখারি ৫৭৮৪; সহীহ মুসলিম ২০৮৫; সুনানে আবু দাউদ ৪০৯০ \\
\textbf{অধ্যায়} & কিতাবুল লিবাস — বাবু মা যুকিরা ফি ইজার (লুঙ্গি সংক্রান্ত অধ্যায়) \\
\textbf{সনদের মান} & \textbf{সহীহ} \\
\hline
\end{tabular}
\end{table}

\#\#\#\# ২. পটভূমি (\lat{Context})

হাদিসটি অত্যন্ত চমৎকার একটি সংলাপের সাথে এসেছে। নবীজি (সা.) অহংকারে পোশাক টেনে হাঁটার ব্যাপারে সতর্ক করছেন। তখন আবু বকর (রাঃ) — যিনি খিলাফতের আগে থেকেই নবীজির সাথি — জিজ্ঞেস করলেন তাঁর লুঙ্গির এক প্রান্ত নেমে গেলে কি হবে? এটি দেখায় সাহাবারা কতটা সূক্ষ্মভাবে অহংকার এড়িয়ে চলতেন।

\#\#\#\# ৩. পূর্ণ বর্ণনা (বাংলা, \lat{pure})

\begin{quote}
\textbf{\lat{“}যে ব্যক্তি অহংকারে পোশাক টেনে হাঁটে, কিয়ামতের দিন আল্লাহ তার দিকে তাকাবেন না।\lat{”}} \\
আবু বকর (রাঃ) বললেন: \textbf{\lat{“}হে আল্লাহর রাসূল! আমার লুঙ্গির এক প্রান্ত নেমে যায়, যদি না আমি সাবধান থাকি।\lat{”}} \\
নবীজি (সা.) বললেন: \textbf{\lat{“}তুমি তো তা অহংকারে করছ না।\lat{”}}
\end{quote}

\#\#\#\# ৪. ধাপে ধাপে বিশ্লেষণ

\textbf{ধাপ ১ — শাস্তির ঘোষণা:} অহংকারে পোশাক টেনে হাঁটা — আল্লাহর দৃষ্টি থেকে বঞ্চিত হওয়ার কারণ।

\textbf{ধাপ ২ — আবু বকরের প্রশ্ন:} তাঁর লুঙ্গি হেঁটে গেলে নেমে যায় — অজান্তে এমন হলে কী হবে?

\textbf{ধাপ ৩ — নবীজির আশ্বাস:} সেটি অহংকার নয় — কারণ নিয়ত শুদ্ধ। গুরুত্ব নিয়তে — \textbf{\lat{intention} (নিয়ত)}।

\#\#\#\# ৫. শব্দের ব্যাখ্যা (\lat{Vocabulary})

\begin{table}[h]\centering\small
\begin{tabular}{ll}
শব্দ & অর্থ \\
\hline
\textbf{ইজরাহ (\arab{إزاره})} & নিজের লুঙ্গি \\
\textbf{খুয়ালা (\arab{خيلاء})} & অহংকার — প্রদর্শনের চলন \\
\textbf{জররা} & টেনে হাঁটা — মাটিতে টেনে নিয়ে চলা \\
\hline
\end{tabular}
\end{table}

\#\#\#\# ৬. গুরুত্বপূর্ণ শিক্ষা (\lat{Key} \lat{Lessons})

\begin{enumerate}
\item \textbf{অহংকারের নিয়তই মাপকাঠি:} পোশাকের দাম নয়, \textbf{\lat{motivation} (উদ্দেশ্য)} গুরুত্বপূর্ণ।
\item \textbf{সাহাবাদের সতর্কতা:} আবু বকর (রাঃ)-এর মতো বড় সাহাবিও এমন সূক্ষ্ম বিষয়ে প্রশ্ন করেছেন — আমাদেরও করা উচিত।
\item \textbf{ইসলাম চরমপন্থা নয়:} পরিচ্ছন্ন পোশাক ও শালীনতা সুন্নাত; শুধু অহংকারের ভাব নিষিদ্ধ।
\end{enumerate}
\#\#\#\# ৭. জীবনে প্রয়োগ (\lat{Practical} \lat{Application})

\begin{enumerate}
\item পোশাক পরার সময় নিয়ত করুন — পরিচ্ছন্নতা ও শালীনতার জন্য, প্রদর্শনের জন্য নয়।
\item অজান্তে কোনো \textbf{\lat{mistake} (ভুল)} হলে বেশি চিন্তা করবেন না — নিয়ত শুদ্ধ রাখুন।
\end{enumerate}
\medskip\hrule\medskip

\subsubsection*{হাদিস ৫ — অহংকারে হাঁটা, পৃথিবী গ্রাস (সহীহ বুখারি ৫৭৮৮, মুসলিম ২০৮৮)}

\#\#\#\# ১. সূত্র কার্ড

\begin{table}[h]\centering\small
\begin{tabular}{ll}
বিষয় & বিবরণ \\
\hline
\textbf{বর্ণনাকারী} & আবু হুরাইরা (রাঃ) \\
\textbf{গ্রন্থ} & সহীহ বুখারি ৫৭৮৮; সহীহ মুসলিম ২০৮৮; রিয়াদুস সালিহীন ৬১৮ \\
\textbf{অধ্যায়} & কিতাবুল লিবাস \\
\textbf{সনদের মান} & \textbf{সহীহ} (মুত্তাফাকুন আলাইহি) \\
\hline
\end{tabular}
\end{table}

\#\#\#\# ২. পটভূমি (\lat{Context})

এটি একটি চমকপ্রদ ঘটনার বর্ণনা — বনী ইসরাঈলের এক ব্যক্তি, যে নিজের সৌন্দর্যে ও পোশাকে এমন অহংকারী ছিল যে পৃথিবী তাকে গিলে ফেলল। নবীজি (সা.) এটিকে কিয়ামত পর্যন্ত অব্যাহত \textbf{\lat{punishment} (শাস্তি)} হিসেবে বর্ণনা করেছেন।

\#\#\#\# ৩. পূর্ণ বর্ণনা (বাংলা, \lat{pure})

\begin{quote}
\textbf{\lat{“}এক ব্যক্তি সুন্দর পোশাক পরে, কেশ সাজিয়ে, মাথা হেলিয়ে, অহংকারে হাঁটছিল। আল্লাহ পৃথিবীকে তাকে গ্রাস করার আদেশ দিলেন। সে কিয়ামত পর্যন্ত পৃথিবীর মধ্যে অবতীর্ণ (ডুবতে) থাকবে।\lat{”}}
\end{quote}

\#\#\#\# ৪. ধাপে ধাপে বিশ্লেষণ

\textbf{ধাপ ১ — ব্যক্তির অবস্থা:} সুন্দর পোশাক, সাজানো চুল, অহংকারী চলন — বাহ্যিক \textbf{\lat{appearance} (চেহারা)}-এ মগ্ন।

\textbf{ধাপ ২ — আল্লাহর আদেশ:} পৃথিবী তাকে গ্রাস করার আদেশ পায় — অহংকারের দ্রুত শাস্তি।

\textbf{ধাপ ৩ — স্থায়ী শাস্তি:} কিয়ামত পর্যন্ত ডুবতে থাকা — অহংকারের শাস্তি দীর্ঘস্থায়ী।

\#\#\#\# ৫. শব্দের ব্যাখ্যা (\lat{Vocabulary})

\begin{table}[h]\centering\small
\begin{tabular}{ll}
শব্দ & অর্থ \\
\hline
\textbf{হুল্লাহ (\arab{حلة})} & সুন্দর পোশাক \\
\textbf{মুরাজ্জিল (\arab{مرجل})} & চুল সাজানো — আঁচড়ানো \\
\textbf{খাসাফা} & গ্রাস করা — গিলে ফেলা \\
\textbf{ইয়াতাজালজাল} & ডুবতে থাকা — গভীরে নেমে যাওয়া \\
\hline
\end{tabular}
\end{table}

\#\#\#\# ৬. গুরুত্বপূর্ণ শিক্ষা (\lat{Key} \lat{Lessons})

\begin{enumerate}
\item \textbf{বাহ্যিক সৌন্দর্যে অহংকার:} পোশাক, চুল, চেহারা — এসব নিয়ে গর্ব করা কিবর; নবীজি এই ঘটনা দিয়ে তা স্পষ্ট করেছেন।
\item \textbf{শাস্তি দ্রুত আসতে পারে:} অহংকারের শাস্তি দুনিয়াতেও আসতে পারে — যেমন পৃথিবী গ্রাস।
\item \textbf{দীর্ঘস্থায়ী পরিণতি:} কিয়ামত পর্যন্ত ডুবতে থাকা — অহংকারের পরিণতি সাময়িক নয়।
\end{enumerate}
\#\#\#\# ৭. জীবনে প্রয়োগ (\lat{Practical} \lat{Application})

\begin{enumerate}
\item আয়না দেখে নিজেকে সাজানোর সময় \textbf{\lat{pride} (গর্ব)} না করে শুকরিয়া করুন।
\item অহংকারী চলন-ভঙ্গি এড়িয়ে চলুন — হাঁটার সময় শান্ত ও বিনয়ী থাকুন।
\end{enumerate}
\medskip\hrule\medskip

\subsection*{পার্ট ৩ — কিবরের পরিণাম: জাহান্নাম}

\subsubsection*{হাদিস ৬ — তিন ব্যক্তি যাদের সাথে আল্লাহ কথা বলবেন না (সহীহ মুসলিম ১০৭)}

\#\#\#\# ১. সূত্র কার্ড

\begin{table}[h]\centering\small
\begin{tabular}{ll}
বিষয় & বিবরণ \\
\hline
\textbf{বর্ণনাকারী} & আবু হুরাইরা (রাঃ) \\
\textbf{গ্রন্থ} & সহীহ মুসলিম ১০৭; সহীহ বুখারি ১২৩৭; সুনানে নাসাই \\
\textbf{অধ্যায়} & কিতাবুল ইমান \\
\textbf{সনদের মান} & \textbf{সহীহ} \\
\hline
\end{tabular}
\end{table}

\#\#\#\# ২. পটভূমি (\lat{Context})

নবীজি (সা.) তিন প্রকার মানুষকে কিয়ামতের দিন আল্লাহর ভয়াবহ পরিত্যাগের খবর দিচ্ছেন। তৃতীয়টি আমাদের বিষয়ের সাথে সরাসরি সম্পর্কিত — \textbf{অহংকারী গরিব}।

\#\#\#\# ৩. পূর্ণ বর্ণনা (বাংলা, \lat{pure})

\begin{quote}
\textbf{\lat{“}তিন প্রকার মানুষ আছে যাদের সাথে কিয়ামতের দিন আল্লাহ কথা বলবেন না, তাদের দিকে তাকাবেন না, তাদের পবিত্র করবেন না, এবং তাদের জন্য যন্ত্রণাদায়ক শাস্তি আছে: (১) বৃদ্ধ ব্যভিচারী, (২) মিথ্যুক শাসক, (৩) অহংকারী গরিব।\lat{”}}
\end{quote}

\#\#\#\# ৪. ধাপে ধাপে বিশ্লেষণ

\textbf{ধাপ ১ — তিনটি অবস্থা:} (ক) বৃদ্ধ ব্যভিচারী — যার প্রবৃত্তি এখনো নিয়ন্ত্রণে নেই, (খ) মিথ্যুক শাসক — ক্ষমতার অপব্যবহার, (গ) অহংকারী গরিব — যার কাছে গর্ব করার কিছুই নেই, তবু অহংকার করে।

\textbf{ধাপ ২ — অহংকারী গরিবের বিশেষত্ব:} গরিব মানুষের অহংকার সবচেয়ে নিকৃষ্ট — কারণ তার কাছে গর্ব করার সম্পদও নেই, শুধু ভেতরের উদ্ধততা।

\textbf{ধাপ ৩ — শাস্তির ভয়াবহতা:} আল্লাহর কথা না বলা, দিকে না তাকানো, পবিত্র না করা — আল্লাহর পরিত্যাগের চরম রূপ।

\#\#\#\# ৫. শব্দের ব্যাখ্যা (\lat{Vocabulary})

\begin{table}[h]\centering\small
\begin{tabular}{ll}
শব্দ & অর্থ \\
\hline
\textbf{শায়খু যানিন} & বৃদ্ধ ব্যভিচারী \\
\textbf{মালিকু কাযিব} & মিথ্যুক শাসক \\
\textbf{আ'ইর মুস্তাকবির} & অহংকারী গরিব \\
\textbf{লাইন ওয়াল্লাহ} & যন্ত্রণাদায়ক শাস্তি \\
\hline
\end{tabular}
\end{table}

\#\#\#\# ৬. গুরুত্বপূর্ণ শিক্ষা (\lat{Key} \lat{Lessons})

\begin{enumerate}
\item \textbf{অহংকারের উৎস দারিদ্র্যেও থাকতে পারে:} সম্পদ না থাকলেও মনোভাব অহংকারী হতে পারে — এটা ভেতরের রোগ।
\item \textbf{নিজের অবস্থা নিয়ে সন্তুষ্টি:} গরিব যদি নিজের অবস্থা নিয়ে অহংকার করে, তবে এটি আল্লাহর \textbf{\lat{displeasure} (অপছন্দ)}।
\item \textbf{তিনটি পরিত্যাগ:} আল্লাহর কথা, দৃষ্টি ও পবিত্রতা — তিনটি থেকে বঞ্চিত হওয়া।
\end{enumerate}
\#\#\#\# ৭. জীবনে প্রয়োগ (\lat{Practical} \lat{Application})

\begin{enumerate}
\item নিজের আর্থিক অবস্থা যাই হোক — অহংকার নয়, \textbf{\lat{gratitude} (কৃতজ্ঞতা)} করুন।
\item শাসক বা বস হলে সততা বজায় রাখুন — মিথ্যা থেকে বাঁচুন।
\end{enumerate}
\medskip\hrule\medskip

\subsubsection*{হাদিস ৭ — জান্নাতী ও জাহান্নামীর চিহ্ন (সহীহ বুখারি ৬০৭১, মুসলিম ২৮৫৩)}

\#\#\#\# ১. সূত্র কার্ড

\begin{table}[h]\centering\small
\begin{tabular}{ll}
বিষয় & বিবরণ \\
\hline
\textbf{বর্ণনাকারী} & হারিসা ইবনে ওয়াহব (রাঃ) \\
\textbf{গ্রন্থ} & সহীহ বুখারি ৬০৭১; সহীহ মুসলিম ২৮৫৩ \\
\textbf{অধ্যায়} & কিতাবুল আদাব — বাবু কানন নবী \\
\textbf{সনদের মান} & \textbf{সহীহ} \\
\hline
\end{tabular}
\end{table}

\#\#\#\# ২. পটভূমি (\lat{Context})

নবীজি (সা.) দুটি সম্প্রদায়ের চারিত্রিক পরিচয় দিচ্ছেন — জান্নাতবাসী ও জাহান্নামবাসী। এই পরিচয় আমাদের নিজেকে যাচাই করার \textbf{\lat{scale} (মানদণ্ড)} দেয়।

\#\#\#\# ৩. পূর্ণ বর্ণনা (বাংলা, \lat{pure})

\begin{quote}
\textbf{\lat{“}কি আমি তোমাদের জান্নাতের অধিবাসীদের সম্পর্কে খবর দেব? সবই দুর্বল, উপেক্ষিত, বিনীত মানুষ। আর কি আমি জাহান্নামের অধিবাসীদের সম্পর্কে খবর দেব? সবই উদ্ধত, কঠোর, অহংকারী।\lat{”}}
\end{quote}

\#\#\#\# ৪. ধাপে ধাপে বিশ্লেষণ

\textbf{ধাপ ১ — জান্নাতীদের চরিত্র:} দুর্বল (শারীরিকভাবে নয়, দুনিয়ার দৃষ্টিতে), উপেক্ষিত, বিনীত — যারা নিজেকে বড় ভাবে না।

\textbf{ধাপ ২ — জাহান্নামীদের চরিত্র:} উদ্ধত, কঠোর (অন্যের প্রতি), অহংকারী — যারা মানুষের সাথে নিষ্ঠুর।

\textbf{ধাপ ৩ — তুলনা:} বিনয় জান্নাতের চাবি, অহংকার জাহান্নামের পথ।

\#\#\#\# ৫. শব্দের ব্যাখ্যা (\lat{Vocabulary})

\begin{table}[h]\centering\small
\begin{tabular}{ll}
শব্দ & অর্থ \\
\hline
\textbf{দু'আফা (\arab{ضعفاء})} & দুর্বল — দুনিয়ার দৃষ্টিতে অক্ষম \\
\textbf{মুস্তাযাফিন} & উপেক্ষিত — সামাজিকভাবে ছোট \\
\textbf{গিলায (\arab{غلاظ})} & কঠোর — রূঢ় \\
\textbf{আতিল (\arab{عتل})} & উদ্ধত — জেদি \\
\textbf{জাওয়াজিয (\arab{جوارع})} & অহংকারী — দাম্ভিক \\
\hline
\end{tabular}
\end{table}

\#\#\#\# ৬. গুরুত্বপূর্ণ শিক্ষা (\lat{Key} \lat{Lessons})

\begin{enumerate}
\item \textbf{বিনয় জান্নাতের চিহ্ন:} যারা দুনিয়ায় নিজেকে বড় ভাবে না, তারা আখেরাতে বড় হবে।
\item \textbf{অহংকার জাহান্নামের চিহ্ন:} উদ্ধত ও কঠোর চরিত্র জাহান্নামী চরিত্র।
\item \textbf{চরিত্র যাচাই:} এই হাদিস দিয়ে নিজের চরিত্র পরীক্ষা করুন — আমি কোন দলের \textbf{\lat{character} (চরিত্র)} ধারণ করছি?
\end{enumerate}
\#\#\#\# ৭. জীবনে প্রয়োগ (\lat{Practical} \lat{Application})

\begin{enumerate}
\item মানুষের সাথে নরম ও \textbf{\lat{gentle} (মৃদু)} থাকুন — কঠোরতা নয়।
\item নিজেকে সামাজিকভাবে ছোট মনে করলে হতাশ নয় — এটিই জান্নাতী চরিত্রের দিকে।
\end{enumerate}
\medskip\hrule\medskip

\subsubsection*{হাদিস ৮ — প্রদর্শনেচ্ছায় আমল: প্রথম তিনজন (সহীহ মুসলিম ১৯০৫)}

\#\#\#\# ১. সূত্র কার্ড

\begin{table}[h]\centering\small
\begin{tabular}{ll}
বিষয় & বিবরণ \\
\hline
\textbf{বর্ণনাকারী} & আবু হুরাইরা (রাঃ) \\
\textbf{গ্রন্থ} & সহীহ মুসলিম ১৯০৫; সুনানে নাসাই; সুনানে তিরমিজি ২৩৮২ \\
\textbf{অধ্যায়} & কিতাবুল ইমারা — বাবু মান কাতালা লির রিয়া \\
\textbf{সনদের মান} & \textbf{সহীহ} \\
\hline
\end{tabular}
\end{table}

\#\#\#\# ২. পটভূমি (\lat{Context})

এটি কিয়ামতের ভয়াবহতম ঘটনা — যেখানে মানুষ তার নেক আমল নিয়ে এসে দেখবে, নিয়ত ছিল রিয়া (প্রদর্শন)। উজব ও রিয়ার সাথে এই হাদিসের সম্পর্ক — মানুষ নিজের আমল নিয়ে আত্মতুষ্ট (উজব) ও প্রদর্শন (রিয়া) করলে আমল বিফল।

\#\#\#\# ৩. পূর্ণ বর্ণনা (বাংলা, \lat{pure})

\begin{quote}
আবু হুরাইরা (রাঃ) থেকে: \textbf{\lat{“}কিয়ামতের দিন প্রথম যাদের বিচার করা হবে, তিনি তিনজন:} \\
১. \textbf{শহীদ} — বলবেন: "আপনার পথে লড়ে শহীদ হলাম।" আল্লাহ বলবেন: "তুমি মিথ্যা বলছ; তুমি লড়েছ যাতে লোক তোমাকে বীর বলে।" \\
২. \textbf{আলেম/কারী} — বলবেন: "আপনার জন্য জ্ঞান অর্জন ও কিরাআত করলাম।" আল্লাহ বলবেন: "তুমি মিথ্যা বলছ; তুমি করেছ যাতে লোক তোমাকে আলেম/কারী বলে।" \\
৩. \textbf{ধনী} — বলবেন: "আপনার পছন্দের সব পথে দান করলাম।" আল্লাহ বলবেন: "তুমি মিথ্যা বলছ; তুমি করেছ যাতে লোক তোমাকে উদার বলে।" \\
সবার পরিণতি: মুখে টেনে জাহান্নামে ফেলা হবে।\lat{”}**
\end{quote}

\#\#\#\# ৪. ধাপে ধাপে বিশ্লেষণ

\textbf{ধাপ ১ — তিনটি নেক আমল:} জিহাদ, ইলম, দান — তিনটি মহান আমল।

\textbf{ধাপ ২ — নিয়তের ফাঁক:} সবাই মানুষের প্রশংসা চেয়েছিল — রিয়া। আমলের বাহ্যিক রূপ নেক, কিন্তু নিয়ত নষ্ট।

\textbf{ধাপ ৩ — আল্লাহর জবাব:} তিনজনের প্রত্যেককে বলা হলো — তুমি মিথ্যা বলছ; তুমি মানুষের প্রশংসা চেয়েছিলে।

\textbf{ধাপ ৪ — পরিণতি:} সবাইকে মুখে টেনে জাহান্নামে ফেলা — রিয়া ও উজবের ভয়াবহ পরিণতি।

\#\#\#\# ৫. শব্দের ব্যাখ্যা (\lat{Vocabulary})

\begin{table}[h]\centering\small
\begin{tabular}{ll}
শব্দ & অর্থ \\
\hline
\textbf{কুতিলা} & শহীদ — যুদ্ধে নিহত \\
\textbf{কুরআন} & কারী — কুরআন তিলাওয়াতকারী \\
\textbf{সুখি} & উদার — দানশীল \\
\textbf{রিয়া} & প্রদর্শন — মানুষের দেখানোর জন্য আমল \\
\textbf{উজব} & আত্মতুষ্টি — নিজের আমলকে বড় ভাবা \\
\hline
\end{tabular}
\end{table}

\#\#\#\# ৬. গুরুত্বপূর্ণ শিক্ষা (\lat{Key} \lat{Lessons})

\begin{enumerate}
\item \textbf{আমল কবুলের শর্ত — ইখলাস:} রিয়া ও উজব আমলকে বিফল করে দেয়।
\item \textbf{প্রথমেই বিচার — বিপদ:} নেক আমলের মানুষও এই ফাঁদে পড়ে — সতর্কতা প্রয়োজন।
\item \textbf{আলেম-শহীদেরও পরীক্ষা:} জ্ঞান ও জিহাদ — দুটিই নিয়তের পরীক্ষায় পড়ে।
\end{enumerate}
\#\#\#\# ৭. জীবনে প্রয়োগ (\lat{Practical} \lat{Application})

\begin{enumerate}
\item যেকোনো আমল শুরু করার আগে নিয়ত \textbf{\lat{check} (যাচাই)} করুন — এটা কি আল্লাহর জন্য নাকি মানুষের প্রশংসার জন্য?
\item আমলের পর মানুষের প্রশংসা কামনা করবেন না — \textbf{\lat{sincerity} (ইখলাস)} রক্ষা করুন।
\end{enumerate}
\medskip\hrule\medskip

\subsubsection*{হাদিস ৯ — জাহান্নাম ও জান্নাতের বিবাদ (সহীহ মুসলিম ২৮৬৫)}

\#\#\#\# ১. সূত্র কার্ড

\begin{table}[h]\centering\small
\begin{tabular}{ll}
বিষয় & বিবরণ \\
\hline
\textbf{বর্ণনাকারী} & আবু হুরাইরা (রাঃ) \\
\textbf{গ্রন্থ} & সহীহ মুসলিম ২৮৬৫; সহীহ বুখারি ৪৮৫০ \\
\textbf{অধ্যায়} & কিতাবুল জান্নাহ \\
\textbf{সনদের মান} & \textbf{সহীহ} \\
\hline
\end{tabular}
\end{table}

\#\#\#\# ২. পটভূমি (\lat{Context})

জান্নাত ও জাহান্নামের মধ্যে \textbf{\lat{dialogue} (সংলাপ)} — যেখানে প্রত্যেকটি নিজের বাসিন্দার পরিচয় বলে। জাহান্নাম বলল — আমার মধ্যে অহংকারী ও উদ্ধতরা; জান্নাত বলল — আমার মধ্যে দুর্বল ও বিনীতরা।

\#\#\#\# ৩. পূর্ণ বর্ণনা (বাংলা, \lat{pure})

\begin{quote}
\textbf{\lat{“}জাহান্নাম ও জান্নাত পরস্পর বিতর্ক করল। জাহান্নাম বলল: আমার মধ্যে অহংকারী ও উদ্ধতরা আসবে। জান্নাত বলল: আমার মধ্যে দুর্বল ও বিনীতরা আসবে। আল্লাহু তাআলা বললেন: "তুমি (জান্নাত) আমার রহমত, তোমার মাধ্যমে আমি যাকে চাই রহমত করি; আর তুমি (জাহান্নাম) আমার শাস্তি, তোমার মাধ্যমে আমি যাকে চাই শাস্তি দিই। প্রত্যেকের (তোমাদের) আমি পূর্ণ করব।"\lat{”}}
\end{quote}

\#\#\#\# ৪. ধাপে ধাপে বিশ্লেষণ

\textbf{ধাপ ১ — জাহান্নামের পরিচয়:} অহংকারী ও উদ্ধত — যারা হকের সামনে মাথা নত করে না।

\textbf{ধাপ ২ — জান্নাতের পরিচয়:} দুর্বল ও বিনীত — যারা আল্লাহর সামনে নত।

\textbf{ধাপ ৩ — আল্লাহর ঘোষণা:} জান্নাত রহমত, জাহান্নাম শাস্তি — প্রত্যেক পূর্ণ হবে। এর মানে — উভয় সম্প্রদায় নিশ্চিত।

\#\#\#\# ৫. শব্দের ব্যাখ্যা (\lat{Vocabulary})

\begin{table}[h]\centering\small
\begin{tabular}{ll}
শব্দ & অর্থ \\
\hline
\textbf{জব্বারিন} & উদ্ধত — জবরদস্ত \\
\textbf{মুতাকাব্বিরিন} & অহংকারী \\
\textbf{দু'আফা} & দুর্বল \\
\textbf{মাসাকিন} & বিনীত — দীনহীন \\
\hline
\end{tabular}
\end{table}

\#\#\#\# ৬. গুরুত্বপূর্ণ শিক্ষা (\lat{Key} \lat{Lessons})

\begin{enumerate}
\item \textbf{চরিত্রই গন্তব্য ঠিক করে:} অহংকার জাহান্নামের, বিনয় জান্নাতের — এটি আল্লাহর নির্ধারিত পদ্ধতি।
\item \textbf{বিনয়ের মহত্ব:} দুনিয়ার দৃষ্টিতে দুর্বল-বিনীতরা আখেরাতে জান্নাতের বাসিন্দা।
\item \textbf{আল্লাহর হিকমত:} জান্নাত রহমতের, জাহান্নাম শাস্তির — আল্লাহ যাকে চান।
\end{enumerate}
\#\#\#\# ৭. জীবনে প্রয়োগ (\lat{Practical} \lat{Application})

\begin{enumerate}
\item বিনয়ী চরিত্র গড়ার জন্য সুন্নাতের আমল (নম্রতা, ক্ষমা, দান) নিয়মিত করুন।
\item অহংকার অনুভব করলে জান্নাত-জাহান্নামের এই বর্ণনা স্মরণ করুন।
\end{enumerate}
\medskip\hrule\medskip

\subsection*{পার্ট ৪ — বিনয়ের ফজিলত}

\subsubsection*{হাদিস ১০ — বিনয়ের ফলে উচ্চতা (সুনানে ইবনে মাজাহ ৪১৭৬, হাকিম ২০৭)}

\#\#\#\# ১. সূত্র কার্ড

\begin{table}[h]\centering\small
\begin{tabular}{ll}
বিষয় & বিবরণ \\
\hline
\textbf{বর্ণনাকারী} & ইবরাহিম ইবনে আদহামের সূত্রে মাওযুফ/মারফু রেওয়ায়াত (সতর্কতা: শুধু সহীহ শাহিদ সহ) \\
\textbf{গ্রন্থ} & সুনানে ইবনে মাজাহ ৪১৭৬; মুসনাদ আহমদ; আল-হাকিম আল-মুসতাদরাক ২০৭ \\
\textbf{সনদের মান} & মূল সনদে দুর্বলতা থাকলেও \textbf{শাহিদ ও ইজমা-র মাধ্যমে হাসান} — ইমাম হাকিম ও জাহাবি একে গ্রহণ করেছেন; শব্দ: \lat{“}\arab{من} \arab{تواضع} \arab{لله} \arab{رفعه} \arab{الله}\lat{”} \\
\hline
\end{tabular}
\end{table}

\#\#\#\# ২. পটভূমি (\lat{Context})

বিনয় এমন একটি গুণ, যা মানুষকে দুনিয়ার দৃষ্টিতে ছোট দেখায়, কিন্তু আল্লাহর দৃষ্টিতে উচ্চ করে। এই হাদিসটি ইসলামের \textbf{\lat{humility} (বিনয়)}-এর কেন্দ্রীয় ধারণা তুলে ধরে — যত বেশি বিনয়ী, তত বেশি উত্তম।

\#\#\#\# ৩. পূর্ণ বর্ণনা (বাংলা, \lat{pure})

\begin{quote}
\textbf{\lat{“}যে আল্লাহর জন্য বিনয়ী হলো, আল্লাহ তাকে উচ্চ করেন।\lat{”}} \\
(বর্ণনাকারী বলেন:) আর যে দুনিয়ার জন্য নিজেকে বড় ভাবল, আল্লাহ তাকে নিচু করলেন।
\end{quote}

\#\#\#\# ৪. ধাপে ধাপে বিশ্লেষণ

\textbf{ধাপ ১ — শর্ত:} \lat{“}আল্লাহর জন্য\lat{”} — বিনয়ও নিয়তের সাথে আবদ্ধ; মানুষের ভয়ে নয়, আল্লাহর জন্য।

\textbf{ধাপ ২ — ফল:} আল্লাহর উচ্চতা — দুনিয়া ও আখেরাত দুটোতেই; মানুষের চোখে ছোট হওয়া আখেরাতে সম্মানের পথ।

\textbf{ধাপ ৩ — বিপরীত:} নিজেকে বড় ভাবলে আল্লাহ নিচু করেন — যেভাবে শয়তান নিজেকে আদম (আঃ) থেকে বড় ভেবে নিচু হলো।

\#\#\#\# ৫. শব্দের ব্যাখ্যা (\lat{Vocabulary})

\begin{table}[h]\centering\small
\begin{tabular}{ll}
শব্দ & অর্থ \\
\hline
\textbf{তাওয়াদা} & বিনয়ী হওয়া \\
\textbf{রাফাআহু} & উচ্চ করা \\
\textbf{জাবারত} & জবরদস্তি — অহংকার \\
\textbf{জালতাহু} & নিচু করা \\
\hline
\end{tabular}
\end{table}

\#\#\#\# ৬. গুরুত্বপূর্ণ শিক্ষা (\lat{Key} \lat{Lessons})

\begin{enumerate}
\item \textbf{বিনয় উচ্চতার পথ:} আল্লাহর কাছে উচ্চতা নম্রতায় — দুনিয়ার পদমর্যাদায় নয়।
\item \textbf{দুনিয়ার অহংকার আখেরাতে নিচুতা:} জবরদস্তি আখেরাতে অপমানের কারণ।
\item \textbf{উদাহরণ:} শয়তান, ফিরাউন, কারুন — প্রত্যেকে নিজেকে বড় ভেবেছে, আল্লাহ তাদের নিচু করেছেন।
\end{enumerate}
\#\#\#\# ৭. জীবনে প্রয়োগ (\lat{Practical} \lat{Application})

\begin{enumerate}
\item যেকোনো \textbf{\lat{success} (সাফল্য)}-এর পর নিজেকে বড় মনে করবেন না — কৃতিত্ব আল্লাহর।
\item ছোট মানুষদের সাথে নম্র ব্যবহার করুন — এতে বিনয় জমা হয়।
\end{enumerate}
\medskip\hrule\medskip

\subsubsection*{হাদিস ১১ — বিনয়ের জন্য এক মর্যাদা (সহীহ মুসলিম ২৫৮৮, তিরমিজি ২০২৯)}

\#\#\#\# ১. সূত্র কার্ড

\begin{table}[h]\centering\small
\begin{tabular}{ll}
বিষয় & বিবরণ \\
\hline
\textbf{বর্ণনাকারী} & আবু হুরাইরা (রাঃ) \\
\textbf{গ্রন্থ} & সহীহ মুসলিম ২৫৮৮; সুনানে তিরমিজি ২০২৯ \\
\textbf{অধ্যায়} & কিতাবুল বিরর — বাবু মা জায়া ফিল তাওয়াদু \\
\textbf{সনদের মান} & \textbf{সহীহ} \\
\hline
\end{tabular}
\end{table}

\#\#\#\# ২. পটভূমি (\lat{Context})

নবীজি (সা.) সাদাকা ও বিনয়ের ফজিলত একসাথে বর্ণনা করছেন। এটি দেখায় — বিনয় শুধু নৈতিকতা নয়, এটি এমন একটি আমল যার জন্য আল্লাহর পক্ষ থেকে \textbf{\lat{reward} (পুরস্কার)} নির্ধারিত।

\#\#\#\# ৩. পূর্ণ বর্ণনা (বাংলা, \lat{pure})

\begin{quote}
\textbf{\lat{“}কোনো দাসী বান্দার সাদাকা দিলে এবং আল্লাহর জন্য বিনয়ী হলে, আল্লাহ তার মর্যাদা বাড়িয়ে দেন।\lat{”}}
\end{quote}

\#\#\#\# ৪. ধাপে ধাপে বিশ্লেষণ

\textbf{ধাপ ১ — সাদাকা:} দান বরকত কমায় না — বাড়ায়।

\textbf{ধাপ ২ — বিনয়:} আল্লাহর জন্য নম্রতা মর্যাদা বাড়ায় — কমায় না।

\textbf{ধাপ ৩ — নীতিমালা:} ইসলামে প্রতিটি নেক আমল মর্যাদা বাড়ায়; দুনিয়ার দৃষ্টিভঙ্গি (দান-বিনয় কমায়) ভুল।

\#\#\#\# ৫. শব্দের ব্যাখ্যা (\lat{Vocabulary})

\begin{table}[h]\centering\small
\begin{tabular}{ll}
শব্দ & অর্থ \\
\hline
\textbf{সাদাকা} & দান \\
\textbf{তাওয়াদা} & বিনয় \\
\textbf{যিয়াদাত} & বৃদ্ধি — মর্যাদার \\
\hline
\end{tabular}
\end{table}

\#\#\#\# ৬. গুরুত্বপূর্ণ শিক্ষা (\lat{Key} \lat{Lessons})

\begin{enumerate}
\item \textbf{বিনয় ও দান — মর্যাদার উৎস:} দুইটিই আল্লাহর নিকট মর্যাদা বাড়ায়।
\item \textbf{দুনিয়ার হিসাব নয়:} মানুষ মনে করে দান কমায়, বিনয় দুর্বলতা — আসলে উল্টো।
\item \textbf{উজবের বিপরীত:} যেখানে উজব নিজেকে বড় ভাবে, সেখানে বিনয় আল্লাহর কাছে উচ্চ করে।
\end{enumerate}
\#\#\#\# ৭. জীবনে প্রয়োগ (\lat{Practical} \lat{Application})

\begin{enumerate}
\item দানে কার্পণ্য নয় — দানই \textbf{\lat{increase} (বৃদ্ধি)}।
\item মানুষের কাছে ক্ষমা চাওয়া, ভুল স্বীকার করা — এতে বিনয় গড়ে ওঠে।
\end{enumerate}
\medskip\hrule\medskip

\subsection*{পার্ট ৫ — উজব: আত্মদাম্ভিকতা ও আত্মতুষ্টি}

\begin{quote}
\textbf{উজব কী?} উজব হলো নিজের আমল-গুণে আত্মতুষ্টি — নিজের ভালো কাজের কারণে নিজেকে বড় ভাবা ও আল্লাহকে ভুলে যাওয়া। কিবর যখন মানুষকে ছোট দেখে, উজব নিজেকেই বড় দেখে। দুটোই আত্মার রোগ।
\end{quote}

\subsubsection*{হাদিস ১২ — পাপ না থাকলে আত্মতুষ্টির আশঙ্কা (আল-বাজ্জার, আল-বায়হাকি)}

\#\#\#\# ১. সূত্র কার্ড

\begin{table}[h]\centering\small
\begin{tabular}{ll}
বিষয় & বিবরণ \\
\hline
\textbf{বর্ণনাকারী} & আনাস ইবনে মালিক (রাঃ) থেকে — ইমাম বাজ্জার ও বায়হাকি \\
\textbf{গ্রন্থ} & মুসনাদুল বাজ্জার ৬৯৩৬; শুআবুল ঈমান; সিলসিলা সাহীহা ৬৫৮ (আলবানি) \\
\textbf{সনদের মান} & \textbf{হাসান লিগাইরিহি} — আলবানি (সিলসিলা সাহীহা ৬৫৮) অনুযায়ী \\
\hline
\end{tabular}
\end{table}

\#\#\#\# ২. পটভূমি (\lat{Context})

এটি একটি আশ্চর্যজনক সতর্কতা — নবীজি (সা.) বলছেন, যদি মানুষ পাপ না করত, তাহলে আল্লাহ এমন মানুষ আনতেন যারা পাপ করে তাওবা করে। কারণ পাপের তাওবা বিনয় নিয়ে আসে; কিন্তু পাপহীন মানুষ আত্মতুষ্ট (উজব) হয়ে পড়ার আশঙ্কা থাকে।

\#\#\#\# ৩. পূর্ণ বর্ণনা (বাংলা, \lat{pure})

\begin{quote}
\textbf{\lat{“}যদি তোমরা পাপ না করতে, আমি আশঙ্কা করতাম তোমাদের ওপর এমন জিনিস আসবে যা পাপের চেয়েও ভয়ংকর — আত্মতুষ্টি (উজব)।\lat{”}}
\end{quote}

\#\#\#\# ৪. ধাপে ধাপে বিশ্লেষণ

\textbf{ধাপ ১ — পাপহীন অবস্থা:} মানুষ পাপ না করলে নিজেকে বড় ভাববে — \lat{“}আমি তো নিষ্পাপ\lat{”} — এই ভাব উজব।

\textbf{ধাপ ২ — পাপের ভূমিকা:} পাপ মানুষকে অহংকার থেকে রক্ষা করে — তাওবার মাধ্যমে বিনয় আসে।

\textbf{ধাপ ৩ — সতর্কতা:} নবীজি উজবকে পাপের চেয়েও ভয়ংকর বলছেন — কারণ পাপ তাওবায় মাফ হয়, কিন্তু উজব আমলই বিফল করে।

\#\#\#\# ৫. শব্দের ব্যাখ্যা (\lat{Vocabulary})

\begin{table}[h]\centering\small
\begin{tabular}{ll}
শব্দ & অর্থ \\
\hline
\textbf{উজব (\arab{عجب})} & আত্মতুষ্টি — নিজের আমল নিয়ে গর্ব \\
\textbf{লাখাশিয়াত} & আশঙ্কা করা \\
\textbf{আশাদ্দু মিনজানযাম্ব} & পাপের চেয়েও ভয়ংকর \\
\hline
\end{tabular}
\end{table}

\#\#\#\# ৬. গুরুত্বপূর্ণ শিক্ষা (\lat{Key} \lat{Lessons})

\begin{enumerate}
\item \textbf{উজব পাপের চেয়েও ভয়ংকর:} পাপ ক্ষমা হয়, কিন্তু উজব আমলকে নষ্ট করে।
\item \textbf{পাপের পরে তাওবা বিনয় আনে:} পাপ মানুষকে অহংকার থেকে রক্ষা করে — তাওবার মাধ্যমে।
\item \textbf{নিষ্পাপতা নিয়ে অহংকার বিপদ:} নিজেকে পাপমুক্ত ভাবা নিজেই একটি বিপদ।
\end{enumerate}
\#\#\#\# ৭. জীবনে প্রয়োগ (\lat{Practical} \lat{Application})

\begin{enumerate}
\item পাপ হলে নিরাশ নয় — সাথে সাথে তাওবা করুন; তাওবা বিনয় নিয়ে আসে।
\item নিজের আমল-ইবাদতের কথা অন্যদের কাছে বড় করে বলবেন না — \textbf{\lat{humility} (বিনয়)} ধরে রাখুন।
\end{enumerate}
\medskip\hrule\medskip

\subsubsection*{হাদিস ১৩ — ধ্বংসের উৎস: দুই জিনিস (ইবনে মাসউদের আথার)}

\#\#\#\# ১. সূত্র কার্ড

\begin{table}[h]\centering\small
\begin{tabular}{ll}
বিষয় & বিবরণ \\
\hline
\textbf{বর্ণনাকারী} & আবদুল্লাহ ইবনে মাসউদ (রাঃ) — সহীহ সনদের আথার \\
\textbf{গ্রন্থ} & মুসনাদুল বাজ্জার; আল-মুজামুল আওসাত (তাবারানি); হিলয়াতুল আউলিয়া ৪/৩৫৪ \\
\textbf{সনদের মান} & সহীহ সনদের \textbf{আসার (আথার)} — ইমাম তাবারানি বর্ণনা করেছেন \\
\hline
\end{tabular}
\end{table}

\#\#\#\# ২. পটভূমি (\lat{Context})

ইবনে মাসউদ (রাঃ) — যিনি নিজের জীবনে বিনয়ের চরম উদাহরণ দেখিয়েছেন (১০ নম্বর গল্প) — তিনি উজবের চরম বিপদ সম্পর্কে সতর্ক করেছেন।

\#\#\#\# ৩. পূর্ণ বর্ণনা (বাংলা, \lat{pure})

\begin{quote}
\textbf{\lat{“}ধ্বংস দুটি জিনিসে: নিরাশা (কুনূত) আর আত্মতুষ্টি (উজব)।\lat{”}} \\
(ব্যাখ্যা:) নিরাশা মানুষকে আল্লাহর রহমত থেকে বিচ্ছিন্ন করে; আত্মতুষ্টি মানুষকে আল্লাহর দিকে ফিরতে দেয় না।
\end{quote}

\#\#\#\# ৪. ধাপে ধাপে বিশ্লেষণ

\textbf{ধাপ ১ — দুটি চরম:} কুনূত — আল্লাহর রহমত থেকে নিরাশা; উজব — নিজের অবস্থা নিয়ে আত্মতুষ্টি।

\textbf{ধাপ ২ — বিপদ:} উভয়ই মানুষকে আল্লাহ থেকে দূরে রাখে — একটিতে হতাশা, অন্যটিতে অহংকার।

\textbf{ধাপ ৩ — মধ্যপথ:} ইসলামের শিক্ষা — আশা ও ভয়ের ভারসাম্য; আশা ছাড়া নিরাশা, ভয় ছাড়া আত্মতুষ্টি।

\#\#\#\# ৫. শব্দের ব্যাখ্যা (\lat{Vocabulary})

\begin{table}[h]\centering\small
\begin{tabular}{ll}
শব্দ & অর্থ \\
\hline
\textbf{কুনূত (\arab{قنوط})} & নিরাশা — রহমতের আশা ত্যাগ \\
\textbf{উজব (\arab{عجب})} & আত্মতুষ্টি — নিজের অবস্থায় গর্ব \\
\textbf{হালাক} & ধ্বংস \\
\hline
\end{tabular}
\end{table}

\#\#\#\# ৬. গুরুত্বপূর্ণ শিক্ষা (\lat{Key} \lat{Lessons})

\begin{enumerate}
\item \textbf{আশা-ভয়ের ভারসাম্য:} মুমিনের হৃদয়ে আশা ও ভয় দুটোই থাকা দরকার।
\item \textbf{উজব — প্রগতির শত্রু:} যে নিজের অবস্থায় সন্তুষ্ট, সে আর উন্নতি করে না।
\item \textbf{নিরাশাও ধ্বংস:} রহমতের আশা হারানোও ধ্বংস — উজবের মতোই।
\end{enumerate}
\#\#\#\# ৭. জীবনে প্রয়োগ (\lat{Practical} \lat{Application})

\begin{enumerate}
\item নিজের \textbf{\lat{progress} (অগ্রগতি)} নিয়ে সন্তুষ্ট না থেকে আরও ভালো আমল চেষ্টা করুন।
\item পাপের পর নিরাশ না হয়ে তাওবা করুন — রহমতের আশা রাখুন।
\end{enumerate}
\medskip\hrule\medskip

\subsubsection*{হাদিস ১৪ — আত্মতুষ্টিতে আমল নষ্ট (উমারের আথার — মুসান্নাফ ইবনে আবি শায়বাহ)}

\#\#\#\# ১. সূত্র কার্ড

\begin{table}[h]\centering\small
\begin{tabular}{ll}
বিষয় & বিবরণ \\
\hline
\textbf{বর্ণনাকারী} & উমার ইবনুল খাত্তাব (রাঃ) — সহীহ সনদের আথার \\
\textbf{গ্রন্থ} & মুসান্নাফ ইবনে আবি শায়বাহ ৩৭১৮০ \\
\textbf{সনদের মান} & \textbf{সহীহ} সনদের আসার \\
\hline
\end{tabular}
\end{table}

\#\#\#\# ২. পটভূমি (\lat{Context})

উমার (রাঃ) — যিনি ক্ষমতার শীর্ষে থেকেও বিনয়ের চরম উদাহরণ (৮, ৯, ১১, ১২ নম্বর গল্প) — উজব সম্পর্কে ভয়ংকর সতর্কতা দিয়েছেন।

\#\#\#\# ৩. পূর্ণ বর্ণনা (বাংলা, \lat{pure})

\begin{quote}
\textbf{\lat{“}আমলকে নষ্ট করার জন্য আমি আত্মতুষ্টি (উজব)-এর চেয়ে ভয়ংকর আর কিছু জানি না — যে মানুষ ভালো আমল করে, তারপর ভাবে: আমি ভালো মানুষ।\lat{”}}
\end{quote}

\#\#\#\# ৪. ধাপে ধাপে বিশ্লেষণ

\textbf{ধাপ ১ — আমল ও উজব:} আমল আল্লাহর দিকে নিয়ে যায়; উজব আমলটিকে ধ্বংস করে।

\textbf{ধাপ ২ — চিহ্ন:} \lat{“}আমি ভালো মানুষ\lat{”} — এই ভাবই উজব; এটা আমলকে নষ্ট করে।

\textbf{ধাপ ৩ — সতর্কতা:} উমার (রাঃ) উজবকে ভয় করেছেন — এটিই তার বিনয়।

\#\#\#\# ৫. শব্দের ব্যাখ্যা (\lat{Vocabulary})

\begin{table}[h]\centering\small
\begin{tabular}{ll}
শব্দ & অর্থ \\
\hline
\textbf{আতালাফ} & নষ্ট করা \\
\textbf{উজব} & আত্মতুষ্টি \\
\textbf{কুনতু ফুলান} & আমি এমন — নিজের প্রশংসার ভাব \\
\hline
\end{tabular}
\end{table}

\#\#\#\# ৬. গুরুত্বপূর্ণ শিক্ষা (\lat{Key} \lat{Lessons})

\begin{enumerate}
\item \textbf{উজব আমল নষ্ট করে:} ভালো আমল + আত্মতুষ্টি = নষ্ট আমল।
\item \textbf{আমল আল্লাহর নেয়ামত:} আমল করতে পেরেছি — এটি আল্লাহর অনুগ্রহ, আমার যোগ্যতা নয়।
\item \textbf{উমারের বিনয়:} ক্ষমতার অধিকারী উমার উজবকে ভয় করেছেন — আমাদের কেমন ভয় করা উচিত?
\end{enumerate}
\#\#\#\# ৭. জীবনে প্রয়োগ (\lat{Practical} \lat{Application})

\begin{enumerate}
\item ভালো আমল শেষে বলুন: \lat{“}আলহামদুলিল্লাহ, আল্লাহ আমাকে তাওফিক দিয়েছেন\lat{”} — নিজের প্রশংসা নয়।
\item নিজের আমলকে বড় মনে করলে সাথে সাথে \textbf{\lat{forgiveness} (ইস্তিগফার)} ও তাওবা করুন।
\end{enumerate}
\medskip\hrule\medskip

\subsubsection*{হাদিস ১৫ — আমল জান্নাতে নিয়ে যায় না (সহীহ মুসলিম ২৮১৬)}

\#\#\#\# ১. সূত্র কার্ড

\begin{table}[h]\centering\small
\begin{tabular}{ll}
বিষয় & বিবরণ \\
\hline
\textbf{বর্ণনাকারী} & আবু হুরাইরা (রাঃ) \\
\textbf{গ্রন্থ} & সহীহ মুসলিম ২৮১৬ \\
\textbf{অধ্যায়} & কিতাবুল জাহান্নাম — বাবু জাননাত আলা আমালিকুম বিল জান্নাহ \\
\textbf{সনদের মান} & \textbf{সহীহ} \\
\hline
\end{tabular}
\end{table}

\#\#\#\# ২. পটভূমি (\lat{Context})

এটি উজব ও আত্মতুষ্টির সবচেয়ে শক্তিশালী নিরাময়। নবীজি (সা.) স্পষ্ট করছেন — আমল মানুষকে জান্নাতে নিয়ে যায় না; আল্লাহর রহমতই নিয়ে যায়। এর অর্থ আমল অকেজো নয় — বরং আমল ছাড়া রহমত পাওয়া কঠিন, কিন্তু আমল নিজে গ্যারান্টি নয়।

\#\#\#\# ৩. পূর্ণ বর্ণনা (বাংলা, \lat{pure})

\begin{quote}
\textbf{\lat{“}তোমাদের কাউকে তার আমল জান্নাতে দেবে না।\lat{”}} সাহাবারা বললেন: \textbf{\lat{“}আপনাকেও না?\lat{”}} নবীজি বললেন: \textbf{\lat{“}আমাকেও না — যদি না আল্লাহ আমাকে তাঁর রহমত দিয়ে ঢেকে দেন।\lat{”}}
\end{quote}

\#\#\#\# ৪. ধাপে ধাপে বিশ্লেষণ

\textbf{ধাপ ১ — ঘোষণা:} আমল জান্নাতের গ্যারান্টি নয় — এটি রহমতের পথ প্রস্তুত করে।

\textbf{ধাপ ২ — সাহাবাদের প্রশ্ন:} আমাকেও না? — নবীজির জবাব: আমাকেও না, রহমত ছাড়া।

\textbf{ধাপ ৩ — উজবের নিরাময়:} এই হাদিস উজবের মূল শিকড় কাটে — নিজের আমলকে বড় ভাবার অবকাশ নেই; সব রহমতে।

\#\#\#\# ৫. শব্দের ব্যাখ্যা (\lat{Vocabulary})

\begin{table}[h]\centering\small
\begin{tabular}{ll}
শব্দ & অর্থ \\
\hline
\textbf{আমল} & কাজ — ইবাদত \\
\textbf{আজইয়াতা} & প্রবেশ করানো \\
\textbf{রহমত} & দয়া — অনুগ্রহ \\
\hline
\end{tabular}
\end{table}

\#\#\#\# ৬. গুরুত্বপূর্ণ শিক্ষা (\lat{Key} \lat{Lessons})

\begin{enumerate}
\item \textbf{আমল না করলে রহমত পাবে না:} আমল রহমত পাওয়ার উপায়, কিন্তু গ্যারান্টি নয়।
\item \textbf{নবীজির বিনয়:} নবীজিও নিজের আমলের ওপর ভরসা করেননি — রহমতের ওপর ভরসা করেছেন।
\item \textbf{উজবের চিকিৎসা:} নিজের আমল নিয়ে গর্ব করলে এই হাদিস স্মরণ করুন।
\end{enumerate}
\#\#\#\# ৭. জীবনে প্রয়োগ (\lat{Practical} \lat{Application})

\begin{enumerate}
\item আমলের পর \textbf{\lat{hope} (আশা)} ও \textbf{\lat{fear} (ভয়)} দুটো রাখুন — জান্নাতের আশা, আমল নষ্ট হওয়ার ভয়।
\item ইবাদতের পর \lat{“}আমি ভালো করলাম\lat{”} না ভেবে — আল্লাহর রহমত চেয়ে করুন।
\end{enumerate}
\medskip\hrule\medskip

\subsection*{পার্ট ৬ — ইলম ও রিয়া: জ্ঞানের অহংকার}

\subsubsection*{হাদিস ১৬ — জ্ঞান অর্জনের নিয়ত (সুনানে আবু দাউদ ৩৬৬৪, ইবনে মাজাহ ২৫২)}

\#\#\#\# ১. সূত্র কার্ড

\begin{table}[h]\centering\small
\begin{tabular}{ll}
বিষয় & বিবরণ \\
\hline
\textbf{বর্ণনাকারী} & আবু হুরাইরা (রাঃ) \\
\textbf{গ্রন্থ} & সুনানে আবু দাউদ ৩৬৬৪; সুনানে ইবনে মাজাহ ২৫২; মুসনাদ আহমদ ৮৪৫৭; সহীহ ইবনে হিব্বান ৭৮; মুসনাদ আবু আওয়ানা; সিলসিলা সাহীহা ৭৬৯ \\
\textbf{সনদের মান} & \textbf{সহীহ} (আলবানি: সহীহ লিগাইরিহি; হাকিম: সহীহ) \\
\hline
\end{tabular}
\end{table}

\#\#\#\# ২. পটভূমি (\lat{Context})

ইলম অর্জন — সর্বোচ্চ ইবাদত; কিন্তু নিয়ত নষ্ট হলে তা সবচেয়ে বড় ক্ষতি। নবীজি (সা.) সতর্ক করছেন — যে জ্ঞান শেখে শুধু মানুষের প্রশংসার জন্য, তার জন্য জান্নাতের \textbf{বিপদ} ঘোষণা করা হয়েছে।

\#\#\#\# ৩. পূর্ণ বর্ণনা (বাংলা, \lat{pure})

\begin{quote}
\textbf{\lat{“}যে ব্যক্তি এমন জ্ঞান অর্জন করল, যা দ্বারা আল্লাহর সন্তুষ্টি কামনা করা হয়, কিন্তু তা শিখল শুধু দুনিয়ার কোনো বস্তু (সম্পদ/মান-সম্মান) পাওয়ার জন্য, সে জান্নাতের সুগন্ধিও পাবে না।\lat{”}}
\end{quote}

\#\#\#\# ৪. ধাপে ধাপে বিশ্লেষণ

\textbf{ধাপ ১ — জ্ঞানের মহত্ব:} ইলম — আল্লাহর সন্তুষ্টির পথ; সবচেয়ে বড় ইবাদত।

\textbf{ধাপ ২ — নিয়তের নষ্ট:} শিখল দুনিয়ার বস্তুর জন্য — সম্পদ, পদ, সম্মান — রিয়া।

\textbf{ধাপ ৩ — শাস্তি:} জান্নাতের সুগন্ধি পাবে না — জ্ঞানের অপব্যবহারের কঠিন শাস্তি।

\#\#\#\# ৫. শব্দের ব্যাখ্যা (\lat{Vocabulary})

\begin{table}[h]\centering\small
\begin{tabular}{ll}
শব্দ & অর্থ \\
\hline
\textbf{ইলম (\arab{علم})} & জ্ঞান \\
\textbf{উজহুহুল্লাহ (\arab{وجه} \arab{الله})} & আল্লাহর সন্তুষ্টি \\
\textbf{আরায (\arab{عرض})} & দুনিয়ার বস্তু — সম্পদ \\
\textbf{রিয়া (\arab{رياء})} & প্রদর্শন — মানুষের দেখানো \\
\hline
\end{tabular}
\end{table}

\#\#\#\# ৬. গুরুত্বপূর্ণ শিক্ষা (\lat{Key} \lat{Lessons})

\begin{enumerate}
\item \textbf{ইলম অর্জনের নিয়ত — অপরিহার্য:} ইলম নিয়তের মাধ্যমে ইবাদত বা বিপদ হয়ে যায়।
\item \textbf{জ্ঞানীর দায়বদ্ধতা:} যত বেশি জ্ঞান, তত বেশি \textbf{\lat{accountability} (দায়বদ্ধতা)}।
\item \textbf{আলেমের অহংকার সবচেয়ে নিকৃষ্ট:} জ্ঞানকে ব্যবহার করা মানুষের প্রশংসার জন্য — উজব ও রিয়া দুটোই।
\end{enumerate}
\#\#\#\# ৭. জীবনে প্রয়োগ (\lat{Practical} \lat{Application})

\begin{enumerate}
\item ইলম শেখার \textbf{\lat{intention} (নিয়ত)} করুন — নিজে জানা, আমল করা ও প্রচার করা।
\item জ্ঞান নিয়ে বড়াই না করে — জ্ঞানের সাথে বিনয় রাখুন।
\end{enumerate}
\medskip\hrule\medskip

\subsubsection*{হাদিস ১৭ — ইলমে বড়াই নিষেধ (সুনানে ইবনে মাজাহ ২৫৪, ২৫৯)}

\#\#\#\# ১. সূত্র কার্ড

\begin{table}[h]\centering\small
\begin{tabular}{ll}
বিষয় & বিবরণ \\
\hline
\textbf{বর্ণনাকারী} & জাবির ইবনে আবদিল্লাহ ও হুযায়ফা (রাঃ) \\
\textbf{গ্রন্থ} & সুনানে ইবনে মাজাহ ২৫৪, ২৫৯; সিলসিলা সাহীহা ২৯৬৪ \\
\textbf{সনদের মান} & \textbf{সহীহ} (আলবানি: সহীহ লিগাইরিহি; শাহিদসহ) \\
\hline
\end{tabular}
\end{table}

\#\#\#\# ২. পটভূমি (\lat{Context})

জাবির (রাঃ) ও হুযায়ফা (রাঃ)-এর দুটি বর্ণনা একই সতর্কতা — আলেমরা যেন জ্ঞানকে দুনিয়ার মুনাফার জন্য না ব্যবহার করে।

\#\#\#\# ৩. পূর্ণ বর্ণনা (বাংলা, \lat{pure})

\begin{quote}
\textbf{\lat{“}জ্ঞান অর্জন করো না আলেমদের সাথে বড়াই করার জন্য, মূর্খদের সাথে ঝগড়া করার জন্য, বা মানুষের দৃষ্টি আকর্ষণ করার জন্য। যে এমন করবে, সে জাহান্নামে যাবে।\lat{”}} \\
(হুযায়ফার বর্ণনায়:) \textbf{\lat{“}যে ব্যক্তি এমন জ্ঞান শেখে যাতে মানুষ তাকে প্রশংসা করে, সে জান্নাতের সুগন্ধিও পাবে না।\lat{”}}
\end{quote}

\#\#\#\# ৪. ধাপে ধাপে বিশ্লেষণ

\textbf{ধাপ ১ — তিনটি নিষেধ:} (ক) আলেমদের সাথে প্রতিযোগিতা (বড়াই), (খ) মূর্খদের সাথে ঝগড়া, (গ) মানুষের দৃষ্টি আকর্ষণ।

\textbf{ধাপ ২ — পরিণতি:} জাহান্নাম — এমন জ্ঞানের ভয়াবহ পরিণতি।

\textbf{ধাপ ৩ — ইলমের ব্যবহার:} ইলমের সঠিক ব্যবহার — আল্লাহর জন্য শেখা, আমল করা, মানুষকে দাওয়াত দেওয়া।

\#\#\#\# ৫. শব্দের ব্যাখ্যা (\lat{Vocabulary})

\begin{table}[h]\centering\small
\begin{tabular}{ll}
শব্দ & অর্থ \\
\hline
\textbf{লি তুবাহি} & বড়াই করার জন্য \\
\textbf{লি তামারি} & ঝগড়া করার জন্য \\
\textbf{লি তাসরিফ} & দৃষ্টি আকর্ষণের জন্য \\
\textbf{আল-উলামা} & আলেমরা \\
\textbf{আস-সুফাহা} & মূর্খরা \\
\hline
\end{tabular}
\end{table}

\#\#\#\# ৬. গুরুত্বপূর্ণ শিক্ষা (\lat{Key} \lat{Lessons})

\begin{enumerate}
\item \textbf{ইলমের নিয়ত পরীক্ষা:} কেন শিখছি? — এটিই প্রতিদিনের প্রশ্ন।
\item \textbf{তিনটি নিষেধ তিনটি বিপদ:} বড়াই, ঝগড়া, দৃষ্টি আকর্ষণ — তিনটি উজবের রূপ।
\item \textbf{ইলমের সাথে বিনয়:} যত জ্ঞান বাড়ে, তত বিনয় বাড়া উচিত।
\end{enumerate}
\#\#\#\# ৭. জীবনে প্রয়োগ (\lat{Practical} \lat{Application})

\begin{enumerate}
\item ইলম শেখার সময় কারো সাথে \textbf{\lat{competition} (প্রতিযোগিতা)} করবেন না — আল্লাহর জন্য শিখুন।
\item জ্ঞানের মজলিসে বিনয়ী শুনুন — বড়াই নয়; প্রয়োজনে নীরব থাকুন।
\end{enumerate}
\medskip\hrule\medskip

\subsection*{সারসংক্ষেপ — এক নজরে}

\begin{table}[h]\centering\small
\begin{tabular}{llll}
বিষয় & হাদিস/আসার & সূত্র & মান \\
\hline
কিবরের সংজ্ঞা ও নিষেধ & সরিষার দানা + সত্য তুচ্ছ ও মানুষ হেয় করা & মুসলিম ৯১ & সহীহ \\
ঈমান ও কিবরের তুলনা & সামান্য ঈমানও জাহান্নাম থেকে রক্ষা & মুসলিম ৯১বি & সহীহ \\
কিবর আল্লাহর চাদর & কিবর-মহিমা আল্লাহর একান্ত গুণ & মুসলিম ২৬২০ & সহীহ \\
পোশাক টানার নিষেধ & অহংকারে পোশাক টানলে আল্লাহ তাকাবেন না & বুখারি ৫৭৮৪ & সহীহ \\
অহংকারে হাঁটা — পৃথিবী গ্রাস & সৌন্দর্যে অহংকারীর শাস্তি & বুখারি ৫৭৮৮ & সহীহ \\
তিন ব্যক্তি আল্লাহর পরিত্যাগে & অহংকারী গরিব & মুসলিম ১০৭ & সহীহ \\
জান্নাতী-জাহান্নামী চরিত্র & বিনীত জান্নাতী, উদ্ধত জাহান্নামী & বুখারি ৬০৭১ & সহীহ \\
রিয়ায় আমল নষ্ট & শহীদ-আলেম-ধনীর বিচার & মুসলিম ১৯০৫ & সহীহ \\
জান্নাত-জাহান্নামের সংলাপ & অহংকারী জাহান্নামে, বিনীত জান্নাতে & মুসলিম ২৮৬৫ & সহীহ \\
বিনয়ের ফজিলত & \lat{“}যে বিনয়ী হলো, আল্লাহ উচ্চ করেন\lat{”} & ইবনে মাজাহ ৪১৭৬ & হাসান \\
বিনয় মর্যাদা বাড়ায় & সাদাকা ও বিনয় & মুসলিম ২৫৮৮ & সহীহ \\
উজবের সতর্কতা & পাপহীন হলে আত্মতুষ্টির আশঙ্কা & সিলসিলা ৬৫৮ & হাসান লিগাইরিহি \\
উজব ধ্বংসের উৎস & নিরাশা + আত্মতুষ্টি & ইবনে মাসউদের আথার & সহীহ \\
উজব আমল নষ্ট করে & \lat{“}আমি ভালো মানুষ\lat{”} ভাব & মুসান্নাফ ৩৭১৮০ & সহীহ \\
আমল জান্নাতে নেয় না & নবীজির বিনয়ের ঘোষণা & মুসলিম ২৮১৬ & সহীহ \\
ইলমের নিয়ত & দুনিয়ার জন্য শেখা — জান্নাতের সুগন্ধি নেই & আবু দাউদ ৩৬৬৪ & সহীহ \\
ইলমে বড়াই নিষেধ & তিনটি নিষেধ — জাহান্নামের সতর্কতা & ইবনে মাজাহ ২৫৪, ২৫৯ & সহীহ \\
\hline
\end{tabular}
\end{table}

\medskip\hrule\medskip

\subsection*{উৎসগ্রন্থ তালিকা}

\begin{itemize}
\item সহীহ বুখারি (ইমাম বুখারি)
\item সহীহ মুসলিম (ইমাম মুসলিম)
\item সুনানে আবু দাউদ (ইমাম আবু দাউদ)
\item সুনানে তিরমিজি (ইমাম তিরমিজি)
\item সুনানে নাসাই (ইমাম নাসাই)
\item সুনানে ইবনে মাজাহ (ইমাম ইবনে মাজাহ)
\item মুসনাদ আহমদ (ইমাম আহমদ বিন হাম্বল)
\item মুসান্নাফ ইবনে আবি শায়বাহ
\item মুসনাদুল বাজ্জার (ইমাম বাজ্জার)
\item শুআবুল ঈমান (ইমাম বায়হাকি)
\item মুস্তাদরাক (ইমাম হাকিম)
\item সিলসিলাতুস সাহীহাহ (শাইখ আলবানি)
\item সহীহ ইবনে হিব্বান
\item রিয়াদুস সালিহীন (ইমাম নববি)
\item শরহে সহীহ মুসলিম (ইমাম নববি)
\end{itemize}
\begin{quote}
\textbf{দ্রষ্টব্য:} যেসব বর্ণনার সনদ দুর্বল (যেমন তাবারানি ১৪৪০০, \lat{“}তিনি ধ্বংসকারী\lat{”}-র ইরাকির সনদ, শুআবুল ঈমানের শয়তান-বর্ণনা) — সেগুলো শর্ত অনুযায়ী সম্পূর্ণ বাদ দেওয়া হয়েছে।
\end{quote}

\end{document}
